In the previous analysis it seems that string theory cannot do better than field theory when the latter does not exist, at least at the perturbative level where one deals with particles.
Moreover when spacetime becomes singular, the string massive modes are not anymore spectators.
Everything seems to suggest that issues with spacetime singularities are hidden into contact terms and interactions with massive states.
This would explain in an intuitive way why the eikonal approach to gravitational scattering works well: the eikonal is indeed concerned with three point massless interactions.
In fact it appears that the classical and quantum scattering on an electromagnetic wave~\cite{Jackiw:1992:ElectromagneticFieldsMassless} or gravitational wave~\cite{tHooft:1987:GravitonDominanceUltrahighenergy} in \bo and \nbo are well behaved.
From this point of view the ACV approach~\cite{Soldate:1987:PartialwaveUnitarityClosedstring,Amati:1987:SuperstringCollisionsPlanckian} may be more sensible, especially when considering massive external states~\cite{Black:2012:HighEnergyString}.
Finally it seems that all issues are related with the Laplacian associated with the space-like subspace with vanishing volume at the singularity.
If there is a discrete zero eigenvalue the theory develops divergences. 


% vim: ft=tex
